\chapter{FGA representability of the Picard scheme}
\label{chap:Picard_FGAPicRepresentability}
% archon:covers AlgebraicJacobian/Picard/FGAPicRepresentability.lean AlgebraicJacobian/Picard/RigidPushforwardGammaBaseChange.lean AlgebraicJacobian/Picard/RigidPushforwardInstance.lean AlgebraicJacobian/Picard/RigidPushforwardP1Topology.lean AlgebraicJacobian/Picard/GaloisDescent/SemilinearAlgebras.lean AlgebraicJacobian/Picard/StableAffineCover.lean

\section{Setup and motivation}
\label{sec:fga_pic_setup}

Throughout this chapter, fix a base field \(k\) and a smooth proper geometrically
integral curve \(C\) over \(k\) (the curve carrying the Jacobian). The relative
Picard presheaf \(\Pic^\sharp_{C/k} : (\Sch/k)^{op} \to \Set\) of
\cref{chap:Picard_LineBundlePullback} has already been packaged as a
set-valued functor and refined to its abelian-group-valued
\'etale-sheafified version \(\Pic^\sharp_{(C/k)\et}\) in
\cref{chap:Picard_RelPicFunctor}. What remains is to produce a scheme
representing it: the \emph{Picard scheme} \(\Pic_{C/k}\). The present chapter
packages that assembly, together with the abelian-group-scheme refinement.

\medskip

\textbf{Two routes to representability.} The Picard scheme of a curve can be
constructed in two ways, and this blueprint records both, because they need
different inputs and only one of them is the route under construction here.

The \emph{quotient route} is Grothendieck's, as presented by Kleiman \S 4: the
Abel map \(\mathrm{Div}_{C/k} \to \Pic^\sharp_{C/k}\) exhibits the Picard functor
as a quotient of a quasi-projective scheme by a smooth proper equivalence
relation, and Altman--Kleiman descent
(\cref{lem:smooth_proper_quotient}) represents the quotient. Its engine is the
Hilbert scheme \(\Hilb_{C/k} = \Quot_{\mathcal{O}_C/C/k}\) of
\cref{chap:Picard_QuotScheme}, whose open subscheme \(\mathrm{Div}_{C/k}\) of
relative effective divisors carries the source of the Abel map. It is the route of
the proof written out under \cref{thm:fga_pic_representability} below.

That attribution needs one qualification, or the reader takes the chapter's headline
theorem for a construction it is not. Neither route is carried out in full here:
\cref{thm:fga_pic_representability} is a statement \emph{about} the existence predicate
of \cref{def:has_pic_scheme}, from which it extracts the representing property, and the
one assertion of this chapter that is not proved is
\cref{thm:fga_pic_representability_etale} --- a bare \texttt{sorry} in Lean, from which
the \(\Pic^\sharp\)-shaped predicate \cref{def:inst_has_pic_scheme} is
\emph{derived} rather than assumed. \Cref{rem:representability_is_conditional} states this
precisely.
So "the route of \cref{thm:fga_pic_representability}" names the argument its proof
records, not a discharge of the predicate --- and the choice of route below is a choice
about which argument is being built out, not about which theorem is available today.

The \emph{Milne--Koll\'ar route} builds \(\Pic^r_{C/k}\) directly, over a
separably closed field, from the open loci
\[
  C^\Sigma = \{\,t : h^0(\mathcal{O}(W - \Sigma)_t) = 1\,\}
\]
cut out inside a degree-\(r\) divisor scheme by a tuple \(\Sigma\) of rational
points, glues them, and descends the result to \(k\) by a finite Galois quotient
before assembling the degrees by a coproduct. Its inputs are a rigidified
pushforward engine, uniform \(H^1\) vanishing, and Galois descent of rigidified
classes --- not a Hilbert or Quot scheme. Its structure, and the parts of it that
are established, are the subject of \cref{sec:fga_pic_milne_kollar}.

The decisive difference is the quotient step. Altman--Kleiman descent needs
quasi-projectivity of the scheme being quotiented, which is not expressible in
the formalisation at present and without which the lemma is false
(\cref{rem:smooth_proper_quotient_hypothesis}); the Milne--Koll\'ar route takes
only finite Galois quotients under an orbit-in-affine hypothesis, and so avoids
that obstruction entirely. The Lean development therefore pursues the
Milne--Koll\'ar route. The quotient-route material below is retained as the
mathematics it is --- the divisor and Hilbert-scheme substrate it rests on is
consumed by both routes --- and is not the path currently being formalised.

\medskip

\textbf{The functor being represented: the \'etale sheafification.} Kleiman \S 4
represents the \'etale-sheafified functor \(\Pic_{(C/k)\et}\), and that is the
formulation this project realises: representability is asserted for
\(\Pic_{(C/k)\et}\) over an \emph{arbitrary} base field, with \emph{no}
hypothesis on \(C(k)\) (\cref{def:pic_et_sheaf},
\cref{thm:fga_pic_representability_etale}).

The reason is not a preference between equivalent options. An unconditional
representability statement for \(\Pic^\sharp_{C/k}\) is not merely unproved but
\emph{false}: Kleiman L5105--L5108 exhibits the conic \(u^2+v^2+w^2=0\) in
\(\I P^2_\bb R\) --- a smooth plane conic, hence smooth, proper and geometrically
integral over a field with no rational point --- for which \(\Pic_{X/\bb R}\) is
not representable while \(\Pic_{(X/\bb R)\et}\) is.

\emph{Corrected} (2026-07-29, \texttt{review-ajc}). This paragraph previously gave
the reason as: \(\Pic^\sharp_{C/k}\) ``is not a sheaf for any topology in general
--- not even the Zariski one'', with representability of a functor forcing the
sheaf condition for every subcanonical topology. The second half is right, but no
source supports the first for the \emph{relative} functor: \texttt{th:cmp} part 1
gives \(\Pic_{X/S}\into\Pic_{(X/S)\zar}\) whenever
\(\mc O_{\!S}\risom f_*\mc O_{\!X}\) holds universally, so on these hypotheses it
is Zariski-\emph{separated}. Take the non-representability directly instead; in
Lean it is
\texttt{PicScheme.not\_exists\_representing\_picSharp\_of\_not\_isIso}, whose only
open input is that the comparison really fails for that conic. Sheafifying repairs this by construction; a rational point repairs
it by hypothesis, and the resulting theorem is about a smaller class of curves.

\medskip

\textbf{What a rational point buys, and why it is kept.} Kleiman \S 2, Thm.~2.5
(\texttt{th:comp}): if \(f : X \to S\) has a section
and \(\mathcal O_{S'} = f'_* \mathcal O_{X'}\) holds for all base changes
(automatic for our proper geometrically integral curve, by cohomology and
base change), then the comparison maps
\[
  \Pic_{X/S}(T) \longrightarrow \Pic_{(X/S)\et}(T) \longrightarrow
  \Pic_{(X/S)\fppf}(T)
\]
are bijective for every \(T\). Hence under a rational-point hypothesis the
plain relative functor \emph{is} the \'etale sheaf, and one may work directly
with \(\Pic^\sharp_{C/k}\).

So there are two formulations, and they prove different theorems.

\begin{enumerate}
  \item \textbf{The one realised.} Represent the \'etale sheafification
    \(\Pic_{(C/k)\et}\) (\cref{def:pic_et_sheaf}) and carry no hypothesis on
    \(C(k)\). This is Kleiman's own formulation, and it needs no section
    precisely because sheafifying supplies \'etale-locally what the section
    would have supplied globally. It is the statement of
    \cref{thm:fga_pic_representability_etale}, and it is what the Jacobian
    headline (\cref{chap:Jacobian}) rests on.
  \item \textbf{Retained, and strictly weaker.} Represent
    \(\Pic^\sharp_{C/k}\) and carry \(C(k) \neq \emptyset\) as a hypothesis.
    This is the formulation of \cref{thm:fga_pic_representability} and of the
    \(\Pic^\sharp\)-shaped assertions later in this chapter. It is strictly
    weaker than a statement about all smooth proper geometrically integral
    curves, since such a curve need not have a rational point --- a conic
    without rational points, or a genus-\(2\) curve over \(\mathbb{Q}\) with no
    rational point, satisfies every other hypothesis. It is retained because it
    is true and because the tangent-space development is written against the
    \(\Pic^\sharp\) shape, and it may not be presented as the project's
    headline.
\end{enumerate}

Which formulation to target was a decision about what theorem is being claimed,
not a limitation of the platform, and it is \emph{settled}: formulation (1), by
owner decision of 2026-07-28. What follows from that decision, and is worth
stating because a reader of the later sections will meet the
\(\Pic^\sharp\)-shaped statements first: those sections' rational-point
hypothesis is now the mark of a conditional milestone rather than a standing
convention of the chapter. A pointed curve \((C,P)\), as in
\cref{chap:AbelJacobi}, supplies the section they need.

% SOURCE: \cite{kleiman-picard}, ``The Picard scheme'', \S 4 (FGA Explained Ch.~9 \S 9.4);
% arXiv:math/0504020, p.~26. Read from
% references/kleiman-picard-src/kleiman-picard.tex.

\section{Line bundles and the Quot / Hilbert scheme}
\label{sec:fga_pic_line_bundle_quot}

The bridge from line bundles on \(C \times_k T\) to the Quot scheme is the
Abel map: a relative effective divisor \(D \subset C \times_k T\) over \(T\)
(a \(T\)-flat closed subscheme whose ideal is invertible) determines a closed
subscheme of \(C \times_k T\) flat over \(T\), hence a \(T\)-point of the Hilbert
functor \(\Hilb_{C/k}(T)\), which is the special case \(\Quot_{\mathcal{O}_C/C/k}\)
of the Quot functor of \cref{chap:Picard_QuotScheme}; and it determines an
invertible sheaf \(\mathcal{O}_{C \times_k T}(D)\) on the relative curve,
representing a class in \(\Pic^\sharp_{C/k}(T)\).

\begin{definition}[Abel-map witness]\leanok
  \label{def:abel_map_witness}
  \dcref{custom}
  \lean{AlgebraicGeometry.Scheme.PicScheme.abelMapWitness}
  \uses{def:div_functor_carrier, def:pic_sharp_carrier}
  Let \(C/k\) be a smooth proper curve. Define a
  natural transformation
  \[
    A^{\mathrm{raw}}_{C/k} \colon \mathrm{Div}_{C/k}
      \longrightarrow \Pic^\sharp_{C/k}
  \]
  as the composite
  \(
    \mathrm{abelKernelNatTrans}(C)
    \mathbin{;} \mathrm{picNeg}(C)
  \).
\end{definition}

\begin{proof}\leanok
  The first transformation sends the class represented by
  \(\langle \mathcal F,q\rangle\) to \([\ker q]\). It is well defined because
  equivalent divisor families induce isomorphic kernels. It is natural under
  arbitrary base change because the kernel--pullback comparison is an
  isomorphism for an epimorphic quotient with finitely presented, base-flat
  target and invertible kernel. The second transformation is pointwise
  inversion in the group-valued presheaf \(\Pic^\sharp_{C/k}\). Their composite
  is therefore natural and sends
  \(\langle \mathcal F,q\rangle\) to \(-[\ker q]\).
\end{proof}

\begin{lemma}[Line bundle / Quot correspondence via the Abel map]\leanok
  \label{lem:line_bundle_quot_correspondence}
  \dcref{custom}
  \lean{AlgebraicGeometry.Scheme.PicScheme.abelMap_app_mk}
  \uses{def:abel_map_witness, def:inst_has_abel_map, def:has_abel_map,
    def:div_functor_carrier, def:pic_sharp_carrier}
  % SOURCE: \cite{kleiman-picard}, ``The Picard scheme'', \S 3, Def.~dfn:Abel (the Abel
  % map), and Thm.~th:repDiv (representability of relative effective
  % divisors by an open subscheme of \(\IHilb_{X/S}\)); arXiv:math/0504020,
  % pp.~25, 23 (read from
  % references/kleiman-picard-src/kleiman-picard.tex, L1837-L1841 + L2135-L2145).
  % Source identifiers: \\label{th:repDiv}\
  % and \\label{dfn:Abel}\.
  \textit{Source: \cite{kleiman-picard}, ``The Picard scheme'', \S 3,
  Thm.~3.7 + Def.~4.6 (cf.\ \cite{nitsure-hilbert-quot}, ``Construction of Hilbert
  and Quot Schemes'', \S 1, where \(\Hilb_{X/S} = \Quot_{\mathcal{O}_X/X/S}\)).}
  Let \(C/k\) be a smooth proper geometrically integral curve. There is a
  natural transformation of presheaves on \((\Sch/k)^{op}\)
  \[
    A_{C/k} = \mathrm{abelMap} \;\colon\; \mathrm{Div}_{C/k}
      \;\longrightarrow\; \Pic^\sharp_{C/k}
  \]
  (\(\mathrm{Div}_{C/k}\) the relative-divisor functor of
  \cref{def:div_functor_carrier}, \(\Pic^\sharp_{C/k}\) the relative
  Picard functor of \cref{def:pic_sharp_carrier}), with the following
  \emph{defining property}: at a \(T\)-point represented by a family
  \(\langle \mathcal F, q\rangle\) (\cref{def:div_family}: a finitely
  presented, \(T\)-flat quotient
  \(q \colon \mathcal O_{C\times_kT}\twoheadrightarrow\mathcal F\)
  whose kernel is an invertible ideal sheaf cutting out a relative effective
  divisor \(D\subseteq C\times_kT\)),
  \[
    A_{C/k}(T)\bigl(\langle \mathcal F, q\rangle\bigr)
      \;=\; -[\ker q] \;=\; [(\ker q)^{-1}]
      \;=\; \bigl[\mathcal{O}_{C \times_k T}(D)\bigr]
      \;\in\; \Pic^\sharp_{C/k}(T),
  \]
  the class of the associated invertible sheaf \(\mathcal{O}_{C \times_k
  T}(D)\), i.e.\ the additive inverse of the class of the ideal sheaf
  \(\ker q\) (Kleiman \S 3, Def.~dfn:Abel). This transformation is the
  required witness for \(\mathrm{HasAbelMap}(C)\)
  (\cref{def:has_abel_map}).
\end{lemma}

\begin{proof}
  The instance of \cref{def:inst_has_abel_map} identifies
  \(A_{C/k}\) with the witness of \cref{def:abel_map_witness}. Evaluating its
  two factors at \(\langle \mathcal F,q\rangle\) first gives the kernel class
  \([\ker q]\) and then its inverse. Since the kernel is the ideal sheaf of
  \(D\), its inverse is \(\mathcal O_{C\times_kT}(D)\), which gives the stated
  formula.
\end{proof}

\begin{remark}
[Representability of \(\mathrm{Div}_{C/k}\)]
  \label{rem:div_functor_representability}
  \dcref{custom}
  Kleiman~\S 3, Thm.~th:repDiv, shows that \(\mathrm{Div}_{C/k}\) is
  representable by an open subscheme \(\mathrm{Div}_{C/k} \subseteq
  \Hilb_{C/k}\) of the Hilbert scheme (\(\Hilb_{C/k} =
  \Quot_{\mathcal{O}_C/C/k}\), \cref{chap:Picard_QuotScheme}): letting
  \(W \subseteq C \times_k \Hilb_{C/k}\) be the universal closed
  subscheme and \(q \colon W \to \Hilb_{C/k}\) the projection, the locus
  \(V \subseteq W\) where \(W\) is locally an effective Cartier divisor
  is open, so \(Z := q(W \setminus V)\) is closed (\(q\) proper) and
  \(\mathrm{Div}_{C/k} := \Hilb_{C/k} \setminus Z\) is the open
  subscheme representing relative effective divisors. This is
  the relative-divisor functor appearing in
  \cref{def:div_functor_carrier}.
\end{remark}

\section{The FGA representability theorem}
\label{sec:fga_pic_representability}

We now state and prove the specialisation of Grothendieck's existence theorem
for the Picard scheme, in the specialisation needed for the Jacobian:
\(S = \Spec k\), \(X = C\) a smooth proper geometrically integral
curve. The general theorem (Kleiman \S 4, Thm.~4.8) covers any
\(X/S\) projective Zariski-locally over \(S\), flat, with integral geometric
fibers; we record the general statement and then specialise.

\begin{theorem}[FGA representability of the Picard scheme]
  \label{thm:fga_pic_representability}
  \dcref{custom}
  \lean{AlgebraicGeometry.Scheme.PicScheme.representable}
  \uses{lem:line_bundle_quot_correspondence, lem:smooth_proper_quotient,
        thm:quot_representable, def:rel_pic_sharp, def:has_rational_point,
        cor:flattening_stratification_curve, def:pic_sharp_representable,
        def:inst_pic_sharp_representable, lem:cech_flat_base_change}
  % SOURCE: \cite{kleiman-picard}, ``The Picard scheme'', \S 4, Main Thm.~th:main,
  % and Cor.~cor:algsch (specialisation to \(S = \Spec k\), \(X\) complete);
  % arXiv:math/0504020, pp.~26, 35 (read from
  % references/kleiman-picard-src/kleiman-picard.tex, L2155-L2166 + L2686-L2690).
  % Source identifiers: \\label{th:main}\
  % and \\label{cor:algsch}\.
  \textit{Source: \cite{kleiman-picard}, ``The Picard scheme'', \S 4, Main Thm.~4.8
  + Cor.~4.18.3.}
  Let \(k\) be a field and \(C\) a smooth proper geometrically integral curve
  over \(k\) \emph{with a \(k\)-rational point}
  (\cref{def:has_rational_point}). Then the relative Picard functor
  \(\Pic^\sharp_{C/k}\) of \cref{def:rel_pic_sharp} --- which under the
  rational-point hypothesis coincides with its \'etale sheafification
  \(\Pic_{(C/k)\et}\), by Kleiman \S 2, Thm.~2.5 --- is representable by a
  \(k\)-scheme \(\Pic_{C/k}\), separated and locally of finite type over
  \(k\), which is a disjoint union of open quasi-projective
  \(k\)-subschemes, each an increasing union of open quasi-projective
  subschemes. The scheme \(\Pic_{C/k}\) is called the \emph{Picard scheme}
  of \(C/k\).

  \emph{Formalisation status} (2026-07-29, \texttt{review-ajc}). The
  \texttt{\textbackslash leanok} mark was removed from this statement: unlike
  \cref{def:inst_pic_sharp_representable}, whose mark truthfully certifies that
  an extraction elaborates, this theorem asserts \emph{representability}, and the
  pinned declaration
  \texttt{PicScheme.representable} does not prove this theorem. It is the
  \texttt{choose\_spec} projection of the hypothesis class
  \texttt{HasPicScheme}, which the declaration itself assumes as an instance
  binder, so what is checked is \(P \to P\) rather than the representability
  above. That class has no instance: its only producer requires a
  \(k\)-rational point (\cref{def:has_rational_point}), and the leaf asserting
  every such curve has one was \emph{false} and has been deleted. The
  mathematical content of this theorem is open, and is carried in Lean by the
  named obligation \texttt{fgaPicardRepresentability}, stated for the
  \'etale sheafification without a rational-point hypothesis
  (\cref{thm:fga_pic_representability_etale}). Note also that the statement
  above --- with its rational-point hypothesis and its
  \(\Pic^\sharp_{C/k}\) target --- is the \emph{conditional} formulation, and
  strictly weaker than what the project claims.
\end{theorem}

\begin{proof}

  We follow Kleiman~\S 4, Thm.~4.8, specialised to \(S = \Spec k\) and
  \(X = C\). The proof proceeds in four conceptual steps, preceded by a
  comparison step; the underlying algebraic engine is the Quot scheme of
  \cref{chap:Picard_QuotScheme} and the relative Picard functor of
  \cref{chap:Picard_RelPicFunctor}, together with the line-bundle / Quot
  correspondence of \cref{lem:line_bundle_quot_correspondence}.

  \emph{Step 0: reduce to the \'etale sheaf via the rational point.}
  By Kleiman \S 2, Thm.~2.5, since \(C \to \Spec k\) has a section
  (\cref{def:has_rational_point}) and \(\mathcal O_T = \pi_{T*}
  \mathcal O_{C \times_k T}\) holds for every \(k\)-scheme \(T\) (cohomology
  and base change for the proper geometrically integral fibers of
  \(C\times_k T/T\)), the comparison map
  \(\Pic^\sharp_{C/k}(T) \to \Pic_{(C/k)\et}(T)\) is bijective for all
  \(T\). It therefore suffices to represent the \'etale sheaf
  \(\Pic_{(C/k)\et}\), which is what Kleiman's theorem produces; the same
  scheme then represents \(\Pic^\sharp_{C/k}\).

  \emph{Step 1: stratify by Hilbert polynomial.}
  For each polynomial \(\phi \in \mathbb{Q}[n]\), let \(P^\phi \subseteq
  \Pic^\sharp_{(C/k)\et}\) be the \'etale subsheaf whose \(T\)-points are
  represented by invertible sheaves \(\mathcal{L}\) on \(C \times_k T\) with
  \(\chi(C_t, \mathcal{L}_t^{-1}(n)) = \phi(n)\) for all \(t \in T\) (the
  Hilbert polynomial relative to the chosen polarisation
  \(\mathcal{O}_C(1)\) on the projective curve \(C\)). The \(P^\phi\) form a
  disjoint open cover of \(\Pic^\sharp_{(C/k)\et}\) (openness via Kleiman's
  flat-base-change argument, \cref{lem:cech_flat_base_change} / Stacks Tag 02KH;
  cf.\ EGA III 7.9.11). It suffices to
  represent each \(P^\phi\) by an increasing union of open quasi-projective
  \(k\)-schemes.

  \emph{Step 2: bound twists by \(m\)-regularity.}
  Given \(\phi\), fix \(m\) large enough that the \(T\)-points of \(P^\phi\) are
  representable by \(\mathcal{L}\) on \(C \times_k T\) satisfying
  \(R^i \pi_{T*} \mathcal{L}(n) = 0\) for all \(i \ge 1\), \(n \ge m\). Twisting
  by \(\mathcal{O}_C(m)\) defines an automorphism of \(\Pic^\sharp_{(C/k)\et}\)
  that carries \(P^\phi\) to \(P^{\phi_0}_0\) for \(\phi_0(n) := \phi(m + n)\).
  Hence it suffices to represent each \(P^{\phi_0}_0\) by a quasi-projective
  \(k\)-scheme.

  \emph{Step 3: factor through the Abel map.}
  Consider the Abel map \(A_{C/k} : \mathrm{Div}_{C/k} \to \Pic^\sharp_{(C/k)\et}\)
  of \cref{lem:line_bundle_quot_correspondence}. The source \(\mathrm{Div}_{C/k}\)
  is representable by an open subscheme of \(\Hilb_{C/k}\) (and hence by an
  open subscheme of the Quot scheme of \cref{thm:quot_representable}), and
  is quasi-projective over \(k\). Pulling back along
  \(P^{\phi_0}_0 \hookrightarrow \Pic^\sharp_{(C/k)\et}\) yields an open
  subscheme \(Z := P^{\phi_0}_0 \times_{\Pic^\sharp_{(C/k)\et}} \mathrm{Div}_{C/k}\)
  of \(\mathrm{Div}_{C/k}\), which is again quasi-projective. The projection
  \(\alpha : Z \to P^{\phi_0}_0\) is a surjection of \'etale sheaves:
  given a \(T\)-point \(\lambda\) of \(P^{\phi_0}_0\) represented by an \'etale
  cover \(T' \to T\) and an invertible sheaf \(\mathcal{L}'\) on \(C \times_k
  T'\) with \(H^1(C_{t'}, \mathcal{L}'_{t'}) = 0\), the projective bundle
  \(\mathbb{P}(\mathcal{Q}) \to T'\) associated to
  \(\mathcal{Q} := \pi_{T'*}\mathcal{L}'\) is smooth over \(T'\); an \'etale
  cover \(T_1 \to T'\) admits a \(T'\)-map \(T_1 \to \mathbb{P}(\mathcal{Q})\)
  (EGA IV 17.16.3 (ii)) whose composition with the universal divisor map
  \(\mathbb{P}(\mathcal{Q}) \to Z\) produces the required \(T_1\)-point of \(Z\).
  Moreover, the first projection \(Z \times_{P^{\phi_0}_0} Z \to Z\) is
  smooth and proper.

  \emph{Step 4: descend by the smooth-proper quotient lemma.}
  Apply Kleiman~\S 4, Lem.~4.9: given a surjection \(\alpha : Z \to P\)
  of \'etale sheaves with \(Z\) quasi-projective, \(R := Z \times_P Z\)
  representable, and the first projection \(R \to Z\) smooth and proper,
  \(P\) is representable by a quasi-projective scheme. In our case
  \(P = P^{\phi_0}_0\), \(Z\) and \(R\) as above, so \(P^{\phi_0}_0\) is
  quasi-projective. Reassembling over \(\phi \in \mathbb{Q}[n]\) and
  over \(m\), the disjoint union of representable \(P^\phi\) pieces represents
  \(\Pic^\sharp_{(C/k)\et}\). The total scheme is separated (as a disjoint
  union of separated quasi-projectives) and locally of finite type. This
  is Kleiman~\S 4, Cor.~4.18.3 in the \(S = \Spec k\),
  \(X = C\) complete case.
\end{proof}

\subsection{The smooth-proper quotient lemma}
\label{subsec:fga_pic_qt}

The structural lemma underlying step 4 above is recorded separately
because it is the abstract content allowing descent from
\(Z = \mathrm{Div}^{\phi_0}_0\) to the representing Picard piece.

\begin{lemma}[Smooth-proper quotient of \'etale sheaves]
  \label{lem:smooth_proper_quotient}
  \dcref{custom}
  \lean{AlgebraicGeometry.Scheme.PicScheme.smoothProperQuotient}
  \uses{def:has_smooth_proper_quotient}
  % \leanok REMOVED 2026-07-29 (review-ajc). The Lean declaration named above
  % does NOT formalise the statement below. It cannot express hypothesis (ii)
  % (quasi-projectivity of Z), so it takes the conclusion as the hypothesis
  % class HasSmoothProperQuotient, whose single field IS P.IsRepresentable:
  % the Lean theorem is P.IsRepresentable -> P.IsRepresentable, with all four
  % numbered hypotheses unused, zero instances and zero call sites. The
  % mathematics below is stated and proved on paper only. See
  % \cref{rem:smooth_proper_quotient_hypothesis} for why (ii) cannot be dropped,
  % and note the Milne--Kollar route needs neither this lemma nor the class.
  % SOURCE: \cite{kleiman-picard}, ``The Picard scheme'', \S 4, Lem.~lm:qt;
  % arXiv:math/0504020, p.~31 (read from
  % references/kleiman-picard-src/kleiman-picard.tex, L2368-L2375).
  % Source identifier: \\label{lm:qt}\.
  \textit{Source: \cite{kleiman-picard}, ``The Picard scheme'', \S 4, Lem.~4.9.}
  Let \(\alpha : Z \to P\) be a morphism of \'etale sheaves on \((\Sch/k)\).
  Set \(R := Z \times_P Z\). Suppose:
  \begin{enumerate}
    \item \(\alpha\) is a surjection of \'etale sheaves;
    \item \(Z\) is representable by a quasi-projective \(k\)-scheme;
    \item \(R\) is representable by a \(k\)-scheme;
    \item the first projection \(R \to Z\) is representable by a smooth
      and proper map.
  \end{enumerate}
  Then \(P\) is representable by a quasi-projective \(k\)-scheme, and \(\alpha\)
  is representable by a smooth map.
\end{lemma}

\begin{proof}
  By (ii) the structure map \(Z \to \Spec k\) is quasi-projective, hence
  separated, so the first projection \(Z \times_k Z \to Z\) is separated.
  By (iv) the first projection \(R \to Z\) is proper; it factors through
  \(h : R \to Z \times_k Z\) over the diagonal embedding (the second
  component is the second projection of \(R \to Z \times_P Z\)). The map
  \(h\) is therefore proper, and a monomorphism (because \(R\) is the graph
  of an equivalence relation on \(Z\), hence injective on \(T\)-points).
  By EGA IV 8.11.5, a proper monomorphism is a closed immersion: \(h\) is a
  closed embedding.

  Hence \(R \hookrightarrow Z \times_k Z\) is a flat (by (iv)) and proper
  equivalence relation on the quasi-projective \(k\)-scheme \(Z\). By the
  theory of flat-and-proper effective equivalence relations on
  quasi-projective schemes (Altman--Kleiman; the construction places
  the graph \(R \subseteq Z \times_k Z\) inside the universal subscheme of
  the Hilbert scheme \(\Hilb_{Z/k}\) and descends to a closed subscheme
  \(Q \subseteq \Hilb_{Z/k}\) representing the quotient), there exist a
  quasi-projective \(k\)-scheme \(Q\) and a faithfully flat, projective map
  \(Z \to Q\) with \(R = Z \times_Q Z\).

  The map \(Z \to Q\) is flat and proper; its fibers, up to field
  extension, coincide with those of \(R \to Z\), which are smooth by
  hypothesis. Hence \(Z \to Q\) is smooth.

  Finally, \(Z \to Q\) is a surjection of \'etale sheaves (an \'etale
  cover of any \(T\)-point of \(Q\) pulls back to a \(T'\)-point of \(Z\), by
  EGA IV 17.16.3 (ii) applied to \(Z \times_Q T \to T\)). In the category
  of \'etale sheaves, both \(Q\) and \(P\) are coequalisers of the pair
  \(R \rightrightarrows Z\); the coequaliser is unique up to unique
  isomorphism, so \(Q\) represents \(P\). Smoothness of \(\alpha : Z \to P\)
  is the smoothness of \(Z \to Q\) just shown.
\end{proof}

\begin{remark}[Quasi-projectivity is load-bearing]
  \label{rem:smooth_proper_quotient_hypothesis}
  \dcref{custom}
  \uses{lem:smooth_proper_quotient}
  Hypothesis (ii) of \cref{lem:smooth_proper_quotient} --- that \(Z\) be
  representable by a \emph{quasi-projective} \(k\)-scheme --- cannot be weakened
  to bare representability. A Hironaka-type free \(\mathbb{Z}/2\)-action on a
  smooth proper non-projective threefold gives a smooth proper equivalence
  relation satisfying (i), (iii) and (iv) whose quotient \'etale sheaf is an
  algebraic space and not a scheme.

  This is what separates the two routes of \cref{sec:fga_pic_setup}. The
  quotient route must supply quasi-projectivity of the Abel-map slice; the
  Milne--Koll\'ar route quotients only by a finite Galois group, under the
  hypothesis that every finite orbit lies in an affine open, which is a
  condition the Hironaka example fails.
\end{remark}

\section{The Milne--Koll\'ar route}
\label{sec:fga_pic_milne_kollar}

This section records the route of \cref{sec:fga_pic_setup} that avoids the
quasi-projectivity obstruction of \cref{rem:smooth_proper_quotient_hypothesis},
in the three parts it decomposes into, and states the two of its inputs that are
established.

Over a separably closed field \(k'\) the construction is Milne \S 4. Fix
\(r \ge 2g+1\) and a degree-\(r\) divisor scheme \(Z\). For a tuple \(\Sigma\) of
\(r-g\) rational points of \(C' = C \times_k k'\), the locus
\[
  C^\Sigma
    = \{\, t \in Z : h^0\bigl(\mathcal{O}(W - \Sigma)_t\bigr) = 1 \,\}
\]
is open, because \(h^0\) is upper semicontinuous and \(h^0 \ge \chi = 1\) in
degree \(g\); adding \(\Sigma\) back to the canonical divisor family of the
rigidified pushforward gives a second map \(C^\Sigma \to Z\), and the equaliser
\(J^\Sigma\) of the two represents the subfunctor of degree-\(r\) classes
\(\mathcal{L}\) with \(h^0(\mathcal{L}(-\Sigma)_t) = 1\) at every point. These
subfunctors cover, so gluing them represents \(\Pic^r_{C'/k'}\). Descent to
\(k\) is a finite Galois quotient, and the degrees assemble as a coproduct over
\(\mathbb{Z}\) through the translation isomorphisms
\(\Pic^d \cong \Pic^r\) given by tensoring with \(\mathcal{O}((r-d) x_0)\).

The two inputs the route rests on that the quotient route does not are the
rigidified pushforward and the finite Galois quotient. The first is local freeness
of \(q_* \mathcal{L}\) together with its compatibility with base change, and it
holds for every curve satisfying the hypotheses of this chapter
(\cref{thm:rigid_pushforward_gate}). The second is Speiser descent together with
the existence of \(\Gamma\)-stable affine neighbourhoods.

\begin{theorem}[Local freeness of the rigidified pushforward]\leanok
  \label{thm:rigid_pushforward_locally_free}
  \lean{AlgebraicGeometry.Adelic.rigidPushforwardLocallyFree_proved}
  \uses{lem:p1_integral}
  \dcref{custom}
  \textit{Source: \cite{mumford-av}, \emph{Abelian Varieties}, II \S 5.}
  Let \(C\) be a smooth proper geometrically integral curve over a field \(k\),
  let \(A\) be a finitely generated \(k\)-algebra and \(q : C_A \to \Spec A\) the
  constant family. Let \(\mathcal{L}\) be an invertible sheaf on \(C_A\) such that
  for every point \(t\) of \(\Spec A\) some two-chart affine cover of the fibre
  \(C_t\) has surjective {\v C}ech difference map on \(\mathcal{L}_t\) --- for a
  separated fibre with affine chart intersection this computes
  \(H^1(C_t, \mathcal{L}_t) = 0\). Then every point of \(\Spec A\) has an open
  neighbourhood on which \(q_* \mathcal{L}\) is free of rank
  \(h^0(C_t, \mathcal{L}_t)\).
\end{theorem}

\begin{remark}[Why the vanishing hypothesis is stated {\v C}ech-locally]
  \label{rem:fiber_h1_vanishing_form}
  \dcref{custom}
  \uses{thm:rigid_pushforward_locally_free}
  The hypothesis of \cref{thm:rigid_pushforward_locally_free} asks for
  \emph{some} two-chart cover of each fibre with surjective difference map,
  rather than for vanishing of the sheaf cohomology \(H^1(C_t, \mathcal{L}_t)\).
  On a separated fibre whose two charts have affine intersection the {\v C}ech
  complex of that cover computes the sheaf cohomology, so the two agree; the
  existential form is the weaker-looking one and is implied by any form of
  \(h^1 = 0\), so a consumer supplying sheaf vanishing supplies this. The
  identification is cover-independent by Mayer--Vietoris, but that
  independence is not itself established here, which is why the hypothesis is
  stated in the form that does not presuppose it.
\end{remark}

\begin{proof}\leanok
  \uses{lem:p1_integral, thm:adelic_exists_finiteMorphismToP1}
  Choose a finite morphism \(\pi : C \to \mathbb{P}^1_k\)
  (\cref{thm:adelic_exists_finiteMorphismToP1}) and base change it to
  \(\pi_A : C_A \to \mathbb{P}^1_A\). Since \(\pi_A\) is finite, it is affine and
  \(\pi_{A*}\) is exact, so \(q_* \mathcal{L} = p_* (\pi_{A*} \mathcal{L})\) for
  \(p : \mathbb{P}^1_A \to \Spec A\), and the fibrewise cohomology of
  \(\pi_{A*}\mathcal{L}\) agrees with that of \(\mathcal{L}\).

  On \(\mathbb{P}^1_A\) the two standard charts form an affine cover with affine
  intersection, so the {\v C}ech complex of any quasi-coherent module is a
  two-term complex \(d\) of \(\Gamma(\Spec A, \mathcal{O})\)-modules with
  \(H^0 = \ker d\) and \(H^1 = \operatorname{coker} d\). Vanishing of \(H^1\) at
  every point makes \(d\) surjective and \(\ker d\) finitely generated
  projective, hence flat, and makes the formation of \(\ker d\) commute with
  arbitrary base change; this is Mumford's two-term finite-free replacement.

  Global sections of the pushforward are \(\ker d\) by the sheaf condition, so
  \(q_* \mathcal{L}\) is finite locally free, and it remains to identify its
  rank. Flatness gives
  \(\operatorname{rank}_t = \dim_{\kappa(t)} \kappa(t) \otimes_A \ker d\), base
  change moves the residue field inside the kernel, and comparing the
  \(\kappa(t)\)-base-changed {\v C}ech square with the {\v C}ech square of the
  fibre identifies that dimension with \(h^0(C_t, \mathcal{L}_t)\). A finite
  locally free module of locally constant rank is free on a neighbourhood of each
  point.
\end{proof}

\begin{lemma}[Global sections of the pushforward commute with base change]\leanok
  \label{lem:gamma_base_change}
  \dcref{custom}
  \lean{AlgebraicGeometry.Adelic.rigidPushforwardGammaBaseChange_proved}
  \uses{thm:rigid_pushforward_locally_free}
  In the situation of \cref{thm:rigid_pushforward_locally_free}, for every ring
  map \(A \to A'\) --- with no finiteness, flatness or noetherian hypothesis on
  \(A'\) --- the canonical map
  \[
    A' \otimes_A \Gamma(C_A, \mathcal{L})
      \longrightarrow \Gamma(C_{A'}, g'^* \mathcal{L})
  \]
  is bijective.
\end{lemma}

\begin{proof}\leanok
  \uses{thm:rigid_pushforward_locally_free}
  Push forward along the finite \(\pi_A : C_A \to \mathbb{P}^1_A\) as in
  \cref{thm:rigid_pushforward_locally_free} and work with the two-term {\v C}ech
  complex \(d\) of the preimage cover. Fibrewise vanishing of \(H^1\) makes \(d\)
  surjective, and for a surjective two-term complex the formation of \(\ker d\)
  is universally exact: base changing the short exact sequence
  \(0 \to \ker d \to K^0 \to K^1 \to 0\) along \(A \to A'\) stays exact because
  \(K^1\) is flat, so \(A' \otimes_A \ker d \to \ker(d \otimes A')\) is
  bijective. Global sections of the pushforward are \(\ker d\) on both sides by
  the sheaf condition, which gives the claim.
\end{proof}

\begin{theorem}[The rigidified pushforward exists for every curve]\leanok
  \label{thm:rigid_pushforward_gate}
  \dcref{custom}
  \lean{AlgebraicGeometry.Adelic.instHasRigidPushforwardOfCurve}
  \uses{thm:rigid_pushforward_locally_free, lem:gamma_base_change}
  Let \(C\) be a smooth proper geometrically integral curve over a field \(k\).
  Then \(C\) admits the rigidified pushforward: for every finitely generated
  \(k\)-algebra \(A\) and every invertible sheaf on \(C_A\) with fibrewise
  vanishing \(H^1\), the pushforward \(q_* \mathcal{L}\) is locally free of rank
  \(h^0\) of the fibre and its formation commutes with arbitrary base change.
\end{theorem}

\begin{proof}\leanok
  \uses{thm:rigid_pushforward_locally_free, lem:gamma_base_change}
  The two conclusions are \cref{thm:rigid_pushforward_locally_free} and, by
  affine-target descent from the module-level statement of
  \cref{lem:gamma_base_change}, the base-change conclusion: over an affine base
  both ends of the adjunction mate are quasi-coherent, and a morphism of
  quasi-coherent modules on an affine scheme is an isomorphism as soon as it is
  bijective on global sections.
\end{proof}

\begin{lemma}[The projective line over a field is integral]\leanok
  \label{lem:p1_integral}
  \dcref{custom}
  \lean{AlgebraicGeometry.Adelic.isIntegral_p1_of_isDomain_charts}
  Let \(k\) be a field. Then \(\mathbb{P}^1_k\) is an integral scheme.
\end{lemma}

\begin{proof}\leanok
  Each standard chart \(V_i \subseteq \mathbb{P}^1_k\) is affine with section
  ring the polynomial ring \(k[T]\) in the chart coordinate, hence an integral
  domain, so each chart is an integral scheme. Reducedness of \(\mathbb{P}^1_k\)
  follows because the two charts cover and reducedness is local. For
  irreducibility, the overlap \(W = V_0 \cap V_1\) is a nonempty open subset of
  the irreducible \(V_0\), hence irreducible and dense in \(V_0\), and
  symmetrically dense in \(V_1\); as \(V_0 \cup V_1 = \mathbb{P}^1_k\), the set
  \(W\) is a dense irreducible subset of \(\mathbb{P}^1_k\), which is therefore
  the closure of an irreducible set and so irreducible.
\end{proof}

\begin{theorem}[Speiser descent]\leanok
  \label{thm:speiser_descent}
  \dcref{custom}
  \lean{AlgebraicJacobian.GaloisDescent.SemilinearAction.descentAlgEquiv}
  Let \(L/K\) be a finite Galois extension with group \(\Gamma\), and let \(A\)
  be an \(L\)-algebra carrying an action of \(\Gamma\) by ring automorphisms
  which is semilinear over the action of \(\Gamma\) on \(L\). Then the canonical
  map
  \[
    L \otimes_K A^\Gamma \longrightarrow A, \qquad a \otimes w \mapsto a w,
  \]
  is an isomorphism of \(L\)-algebras: the invariants form a \(K\)-form of
  \(A\). No finiteness hypothesis on \(A\) is required.
\end{theorem}

\begin{proof}\leanok
  The same statement for semilinear \(\Gamma\)-modules is the classical
  normal-basis argument: writing \(L = K[\Gamma] \cdot \theta\) for a normal
  basis element \(\theta\), the trace \(\sum_{\gamma} \gamma(\theta a) \gamma\)
  exhibits every element of \(A\) as an \(L\)-combination of invariants, and
  \(L\)-independence of the invariant part follows from linear independence of
  the characters \(\gamma\). The map above is the \(L\)-algebra map extending the
  inclusion \(A^\Gamma \hookrightarrow A\), and its underlying additive map is
  the module-level descent map, which the module statement shows bijective.
  Bijectivity of an algebra homomorphism makes it an algebra isomorphism.
\end{proof}

\begin{theorem}[Points of the finite Galois quotient]\leanok
  \label{thm:galois_quotient_points}
  \dcref{custom}
  \lean{AlgebraicJacobian.GaloisDescent.SemilinearAction.invariantAlgHomEquiv}
  \uses{thm:speiser_descent}
  In the situation of \cref{thm:speiser_descent}, for every \(K\)-algebra \(B\)
  the map \(g \mapsto (\mathrm{id}_L \otimes g)\) is a bijection
  \[
    \mathrm{Hom}_{K\text{-alg}}(A^\Gamma, B)
      \;\simeq\;
    \{\, f \in \mathrm{Hom}_{L\text{-alg}}(A, L \otimes_K B) :
       f(\gamma \cdot x) = (\gamma \otimes \mathrm{id})(f(x)) \,\}.
  \]
  Equivalently, for an affine \(K\)-scheme \(T = \Spec B\), the \(T\)-points of
  \(\Spec A^\Gamma\) are the \(\Gamma\)-equivariant \(T_L\)-points of
  \(\Spec A\): the affine scheme \(\Spec A^\Gamma\) is the quotient of
  \(\Spec A\) by \(\Gamma\).
\end{theorem}

\begin{proof}\leanok
  \uses{thm:speiser_descent}
  Given \(g\), transport \(\mathrm{id}_L \otimes g\) through the isomorphism
  \(A \cong L \otimes_K A^\Gamma\) of \cref{thm:speiser_descent}; equivariance is
  the identity \(\gamma \otimes \mathrm{id}\) transported along the same
  isomorphism, which is exactly the compatibility recorded there. Conversely, an
  equivariant \(f\) carries invariants to \(\Gamma\)-fixed elements of
  \(L \otimes_K B\), and for the left-factor action on a base change the fixed
  points are \(1 \otimes B\); so \(f\) restricts to a \(K\)-algebra map
  \(A^\Gamma \to B\). The two constructions are mutually inverse: one direction
  is the computation of the restriction on invariants, and the other holds
  because the invariants span \(A\) over \(L\), so two \(L\)-algebra maps
  agreeing on them agree.
\end{proof}

\begin{lemma}[\(\Gamma\)-stable affine neighbourhoods]\leanok
  \label{lem:stable_affine_open}
  \dcref{custom}
  \lean{AlgebraicJacobian.GaloisDescent.SemilinearGalAction.exists_stable_affineOpen_of_orbitsInAffineOpen}
  \uses{thm:galois_quotient_points}
  Let \(L/K\) be finite Galois with group \(\Gamma\) and let \(X\) be a scheme
  over \(L\) with a semilinear \(\Gamma\)-action all of whose orbits are
  contained in affine opens. Then every point of \(X\) has a \(\Gamma\)-stable
  affine open neighbourhood.
\end{lemma}

\begin{proof}\leanok
  Let \(x \in X\) and let \(U\) be an affine open containing the orbit of \(x\).
  The orbit is finite and lies in the open
  \(W = \bigcap_{\gamma} \gamma^{-1} U\), so by prime avoidance inside the affine
  \(U\) there is a section \(s \in \Gamma(X, U)\) with the whole orbit contained
  in \(D(s) \subseteq W\). Put \(N = \prod_{\gamma} \gamma(s)|_W\). Then
  \(D(N) = \bigcap_{\gamma} \gamma^{-1} D(s)\), which contains \(x\) because the
  orbit lies in \(D(s)\), and is \(\Gamma\)-stable because translating the
  intersection permutes its factors. Finally \(D(N) \subseteq D(s)\) by the
  factor \(\gamma = 1\), so \(D(N)\) is a basic open of the affine \(D(s)\) and
  hence affine.
\end{proof}

\begin{remark}[The orbit hypothesis is load-bearing]
  \label{rem:orbit_in_affine_hypothesis}
  \dcref{custom}
  \uses{lem:stable_affine_open, rem:smooth_proper_quotient_hypothesis}
  The hypothesis of \cref{lem:stable_affine_open} that every orbit lie in an
  affine open cannot be dropped, and it is the same phenomenon as
  \cref{rem:smooth_proper_quotient_hypothesis}: the Hironaka free
  \(\mathbb{Z}/2\)-action on a smooth proper non-projective threefold has an
  orbit contained in no affine open, and its quotient is an algebraic space
  rather than a scheme. The difference between the two routes is not that one
  needs a hypothesis and the other does not, but that this hypothesis is
  available for the schemes the Milne--Koll\'ar route quotients --- glued from
  opens carrying locally closed immersions into a Grassmannian --- whereas
  quasi-projectivity of the Abel-map slice is not currently expressible.
\end{remark}

\section{Group-scheme structure}
\label{sec:fga_pic_group_structure}

The Picard scheme \(\Pic_{C/k}\) inherits an abelian group structure from
the tensor product of invertible sheaves on \(C \times_k T\). The relative
Picard presheaf \(T \mapsto \Pic(C \times_k T)/H_T\) is already an abelian
group at every \(T\) (the quotient of an abelian group by a subgroup),
functorially in \(T\); this lifts to the \'etale sheafification, and then
to a group-scheme structure on the representing scheme by Yoneda.

\begin{theorem}[\(\Pic_{C/k}\) is a \(k\)-group scheme]\leanok
  \label{thm:pic_is_group_scheme}
  \dcref{custom}
  \lean{AlgebraicGeometry.Scheme.PicScheme.groupSchemeStructure}
  \uses{thm:fga_pic_representability, thm:relative_pic_quotient_well_defined,
        def:pic_scheme, def:rel_pic_sharp}

  % SOURCE: \cite{kleiman-picard}, ``The Picard scheme'', \S 2, Def.~df:Pfs (the relative
  % Picard functor takes values in abelian groups) and \S 4 surrounding
  % text (group-scheme structure on \(\IPic_{X/S}\) via tensor product);
  % arXiv:math/0504020, pp.~17, 26 (read from
  % references/kleiman-picard-src/kleiman-picard.tex, L1311-L1318 + L2057-L2069).
  % Source identifiers: \\label{df:Pfs}\ and
  % \\label{df:Psch}\.
  \textit{Source: \cite{kleiman-picard}, ``The Picard scheme'', \S 2, Def.~2.2 +
  \S 4, Def.~4.1 (the Picard scheme).}
  Let \(\Pic_{C/k}\) be the Picard scheme of
  \cref{thm:fga_pic_representability}. Then \(\Pic_{C/k}\) is canonically an
  abelian group scheme over \(k\): the tensor product
  \([\mathcal{L}] + [\mathcal{M}] := [\mathcal{L} \otimes \mathcal{M}]\)
  and inverse \(-[\mathcal{L}] := [\mathcal{L}^{-1}]\) on
  \(\Pic^\sharp_{C/k}(T)\) (well defined because both operations
  preserve isomorphism classes and descend through the subgroup
  \(H_T = \pi_T^*\Pic(T)\) of \cref{thm:relative_pic_quotient_well_defined})
  are natural in \(T\), hence by Yoneda are represented by morphisms
  \(m : \Pic_{C/k} \times_k \Pic_{C/k} \to \Pic_{C/k}\),
  \(i : \Pic_{C/k} \to \Pic_{C/k}\), and
  \(e : \Spec k \to \Pic_{C/k}\) (the trivial line bundle class) satisfying
  the abelian-group axioms. In particular, \(\Pic_{C/k}\) is a \(k\)-group
  scheme locally of finite type.
\end{theorem}

\begin{proof}
\leanok
  The relative Picard functor \(\Pic^\sharp_{C/k}\) of
  \cref{def:rel_pic_sharp} is abelian-group valued, with sum
  \([\mathcal{L}] + [\mathcal{M}] = [\mathcal{L} \otimes \mathcal{M}]\),
  inverse \(-[\mathcal{L}] = [\mathcal{L}^{-1}]\), and unit
  \(0 = [\mathcal{O}_{C \times_k T}]\); these operations descend through
  the quotient by \(H_T = \pi_T^*\Pic(T)\) because \(\pi_T^*\) is itself a
  group homomorphism (\cref{thm:rel_pic_addcommgroup_via_tensorobj}).

  By \cref{thm:fga_pic_representability}, \(\Pic^\sharp_{C/k}\) is
  represented by the scheme \(\Pic_{C/k}\). The Yoneda embedding is fully
  faithful and preserves limits, so a representing object of a
  group-valued presheaf is canonically a group object: the natural
  transformations
  $\Pic^\sharp_{C/k} \times \Pic^\sharp_{C/k} \to \Pic^\sharp_{C/k}$
  (multiplication), \(\Pic^\sharp_{C/k} \to \Pic^\sharp_{C/k}\) (inverse),
  and the unit map \(h_{\Spec k} \to \Pic^\sharp_{C/k}\) pick out unique
  morphisms of representing schemes \(m\), \(i\), \(e\) as displayed, and
  the abelian-group axioms hold at the level of \(T\)-points functorially
  in \(T\), hence on the representing scheme. The local-finite-type
  property is part of \cref{thm:fga_pic_representability}.
\end{proof}

\section{The Picard scheme}
\label{sec:fga_pic_named_definition}

The two preceding theorems determine the following named scheme.

\begin{definition}[The Picard scheme]\leanok
  \label{def:pic_scheme}
  \dcref{custom}
  \lean{AlgebraicGeometry.Scheme.PicScheme}
  \uses{def:has_pic_scheme, def:inst_has_pic_scheme}

  % SOURCE: \cite{kleiman-picard}, ``The Picard scheme'', \S 4, Def.~df:Psch;
  % arXiv:math/0504020, p.~26 (read from
  % references/kleiman-picard-src/kleiman-picard.tex, L2057-L2062).
  % Source identifier: \\label{df:Psch}\.
  \textit{Source: \cite{kleiman-picard}, ``The Picard scheme'', \S 4, Def.~4.1.}
  Let \(C/k\) be a smooth proper geometrically integral curve over a field
  \(k\), with a \(k\)-rational point. The \emph{Picard scheme} \(\Pic_{C/k}\)
  is the \(k\)-scheme representing the relative Picard functor
  \(\Pic^\sharp_{C/k}\) (equivalently, by Kleiman \S 2 Thm.~2.5 under the
  rational point, its \'etale sheafification), supplied by
  \cref{thm:fga_pic_representability}, equipped with the abelian-group-scheme
  structure of \cref{thm:pic_is_group_scheme}. The scheme is separated,
  locally of finite type over \(k\), and a disjoint union of open
  quasi-projective \(k\)-subschemes.
\end{definition}

\section{Construction interfaces}
\label{sec:fga_pic_lean_encoding}

The constructions above fit together through the following interfaces. The
relative Picard presheaf is obtained from the quotient of invertible sheaves
on \(C\times_kT\) by pullbacks from \(T\), and its set-valued version is used
when applying representability. The relative-divisor functor is the
specialisation of the effective Cartier-divisor functor to
\(C\to\operatorname{Spec}k\). A divisor family \(\langle\mathcal F,q\rangle\)
is sent by the Abel map to the inverse class of its kernel:
\[
  \langle\mathcal F,q\rangle\longmapsto-[\ker q].
\]
The FGA theorem applies Hilbert-polynomial stratification, a regularity bound,
the Abel-map factorisation, and the smooth-proper quotient lemma. Tensor
product and dualisation of line bundles then induce addition and inversion on
the representing scheme.

\begin{itemize}
  \item The Picard scheme \cref{def:pic_scheme} is selected by the
    existence assertion \cref{def:has_pic_scheme}.
  \item The functor \cref{def:pic_sharp_carrier} is the set-valued form of
    the relative Picard presheaf of \cref{chap:Picard_RelPicFunctor}.
  \item The Abel map of \cref{lem:line_bundle_quot_correspondence} is the natural
    transformation \(\mathrm{Div}_{C/k} \to \Pic^\sharp_{C/k}\) sending a
    relative divisor family \(\langle \mathcal F, q\rangle\) to
    \(-[\ker q]\), the class of its associated invertible sheaf.
  \item The Altman--Kleiman lemma \cref{lem:smooth_proper_quotient} gives
    descent of a flat-and-proper equivalence relation \(R \subseteq Z
    \times_k Z\) on a quasi-projective \(k\)-scheme \(Z\) to a
    quasi-projective quotient \(Q\) representing the coequalising sheaf.
  \item The theorem \cref{thm:fga_pic_representability} supplies the
    representing scheme.
  \item Yoneda transport of the abelian-group structure on
    \(\Pic^\sharp_{C/k}\) gives the group-scheme structure of
    \cref{thm:pic_is_group_scheme}.
\end{itemize}

The effective-quotient theorem for flat and proper equivalence relations on
quasi-projective schemes is recorded as
\cref{lem:smooth_proper_quotient}. All other interfaces are supplied by the
relative Picard, divisor, and Quot constructions.

\section{Auxiliary existence predicates}
\label{sec:fga_pic_typeclass_scaffolding}

The public constructions use a small collection of existence predicates. Each
predicate records one mathematical object or property needed by the preceding
representability argument, and the associated definitions select the object
when the predicate is inhabited.

\subsection{Basic functors}
\label{subsec:fga_pic_functor_carriers}

\begin{definition}[\(k\)-rational point]\leanok
  \label{def:has_rational_point}
  \dcref{custom}
  \lean{AlgebraicGeometry.Scheme.HasRationalPoint}
  The predicate \(\mathrm{HasRationalPoint}(C)\) asserts that the structural
  morphism \(C \to \Spec k\) admits a section
  \(\sigma : \Spec k \to C\), i.e.\ that \(C\) has a \(k\)-rational point.
  This is the standing hypothesis of the chapter (see
  \cref{sec:fga_pic_setup}): by Kleiman \S 2, Thm.~2.5, it makes the plain
  relative Picard functor \(\Pic^\sharp_{C/k}\) coincide with its \'etale
  (and fppf) sheafification, so representability can be stated against
  \cref{def:rel_pic_sharp} directly. For a pointed curve \((C,P)\), the
  pointing \(P\) supplies this section.
\end{definition}

\begin{proof}\leanok
  This is a definition.
\end{proof}

\begin{definition}[Set-valued relative Picard functor]\leanok
  \label{def:pic_sharp_carrier}
  \dcref{custom}
  \lean{AlgebraicGeometry.Scheme.PicScheme.picSharp}
  \uses{def:rel_pic_sharp}
  The set-valued shadow \(\Pic^\sharp_{C/k} : (\Sch/k)^{op} \to \Set\) of the
  group-valued relative Picard presheaf of \cref{def:rel_pic_sharp}
  (composition with the forgetful functor \(\Ab \to \Set\)).

  A representing object therefore carries natural bijections
  \[
    (T \to \Pic_{C/k}) \simeq
    \Pic(C \times_k T)/\pi_T^*\Pic(T).
  \]
\end{definition}

\begin{proof}\leanok
  \uses{def:rel_pic_sharp}
  Immediate from \cref{def:rel_pic_sharp}: apply the forgetful functor to
  the group-valued presheaf.
\end{proof}

\begin{definition}[Vacuous carrier, not the existence of \(\Div_{C/k}\)]
  \label{def:has_div_functor}
  \dcref{custom}
  \lean{AlgebraicGeometry.Scheme.PicScheme.HasDivFunctor}
  % \leanok REMOVED and statement CORRECTED 2026-07-29 (review-ajc). The former
  % text claimed this predicate asserts that the presheaf category contains the
  % relative-effective-divisor functor of C. It does not: the Lean field is
  %   has_div_functor : Nonempty ((Over (Spec (.of k)))^op \functor Type (u+1))
  % in which C DOES NOT OCCUR. Any functor whatsoever witnesses it, so the
  % predicate is vacuously true and carries no content about C. It was carrying
  % \leanok as "the existence of the relative divisor functor", which is how the
  % vacuity survived. The real functor is \cref{def:div_functor_carrier}.
  The predicate \(\mathrm{HasDivFunctor}(C)\) asserts only that the category of
  presheaves of types on \((\Sch/k)^{op}\) is \emph{nonempty}. The curve \(C\)
  does not occur in the assertion, so it is vacuously true --- a constant functor
  witnesses it --- and it must never be cited as the existence of
  \(\Div_{C/k}\). Use \cref{def:div_functor_carrier} for that; it is a genuine
  definition instantiating \cref{def:div_functor} at \(C \to \Spec k\).
\end{definition}

\begin{proof}
  This is the defining predicate. It is vacuous, as recorded above.
\end{proof}

\begin{definition}[The relative-divisor functor]\leanok
  \label{def:div_functor_carrier}
  \dcref{custom}
  \lean{AlgebraicGeometry.Scheme.PicScheme.divFunctor}
  \uses{def:div_functor}
  The presheaf \(\mathrm{Div}_{C/k} : (\Sch/k)^{op} \to \Set\) is the
  relative-divisor functor \(\Div_{C/k}\) of \cref{def:div_functor},
  instantiated at the structural morphism \(C \to \Spec k\): it sends a
  \(k\)-scheme \(T\) to the set of relative effective Cartier divisors on
  \(C \times_k T\) flat over \(T\).
\end{definition}

\begin{proof}\leanok
  \uses{def:div_functor}
  Immediate from \cref{def:div_functor}: instantiate at
  \(X = C\), \(S = \Spec k\), \(\pi\) the structural morphism.
\end{proof}

\begin{definition}[The vacuous carrier is inhabited]\leanok
  \label{def:inst_has_div_functor}
  \dcref{custom}
  \lean{AlgebraicGeometry.Scheme.PicScheme.instHasDivFunctor}
  \uses{def:has_div_functor, def:div_functor_carrier}
  % Title corrected 2026-07-29 (review-ajc): was "Relative-divisor functor
  % exists", which overstated it. \leanok is retained -- the Lean instance does
  % discharge its (vacuous) statement -- but what it discharges is
  % \cref{def:has_div_functor}, a nonemptiness claim in which C does not occur.
  The relative-divisor functor of \cref{def:div_functor_carrier} is supplied as
  the witness for \cref{def:has_div_functor}. Since that predicate does not
  mention \(C\), any functor would serve equally well: this instance records a
  choice of witness, not the existence of \(\Div_{C/k}\).
\end{definition}

\begin{proof}\leanok
  \uses{def:div_functor_carrier}
  The functor \(\mathrm{Div}_{C/k}\) of \cref{def:div_functor_carrier} is an
  object of the presheaf category, witnessing the (vacuous) existential.
\end{proof}

\subsection{Abel-map predicate}
\label{subsec:fga_pic_abel_map_carrier}

\begin{definition}[Abel-map data]\leanok
  \label{def:has_abel_map}
  \dcref{custom}
  \lean{AlgebraicGeometry.Scheme.PicScheme.HasAbelMap,
    AlgebraicGeometry.Scheme.PicScheme.abelMap}
  \uses{def:div_functor_carrier, def:pic_sharp_carrier}
  The predicate \(\mathrm{HasAbelMap}(C)\) carries a natural
  transformation of presheaves \(\mathrm{abel} \colon \mathrm{Div}_{C/k}
  \to \Pic^\sharp_{C/k}\) between the relative-divisor functor of
  \cref{def:div_functor_carrier} and the relative Picard functor of
  \cref{def:pic_sharp_carrier}; \(\mathrm{abelMap} :=
  \mathrm{HasAbelMap.abel}\) denotes this transformation.
\end{definition}

\begin{proof}
  This is a definition.
\end{proof}

\begin{definition}[Existence of the Abel map]\leanok
  \label{def:inst_has_abel_map}
  \dcref{custom}
  \lean{AlgebraicGeometry.Scheme.PicScheme.instHasAbelMap}
  \uses{def:has_abel_map, def:abel_map_witness}
  For every such curve \(C/k\), the transformation
  \(A^{\mathrm{raw}}_{C/k}\) from \cref{def:abel_map_witness} supplies
  \cref{def:has_abel_map}.
\end{definition}

\begin{proof}
\leanok
  \uses{def:abel_map_witness}
  The natural transformation constructed in \cref{def:abel_map_witness} is a
  term of the required type
  \(\mathrm{Div}_{C/k} \to \Pic^\sharp_{C/k}\), and hence supplies the data of
  \cref{def:has_abel_map}.
\end{proof}

\subsection{Smooth-proper quotient predicate}
\label{subsec:fga_pic_quotient_carrier}

\begin{definition}[Representability of the smooth-proper quotient]\leanok
  \label{def:has_smooth_proper_quotient}
  \dcref{custom}
  \lean{AlgebraicGeometry.Scheme.PicScheme.HasSmoothProperQuotient}
  For a morphism \(\alpha : Z \to P\) of presheaves of types on
  \((\Sch/k)^{op}\), the predicate \(\mathrm{HasSmoothProperQuotient}(\alpha)\)
  asserts that the target \(P\) is representable. It records the conclusion
  of the Altman--Kleiman lemma \cref{lem:smooth_proper_quotient} for the
  quotient occurring in the FGA construction.
\end{definition}

\begin{proof}\leanok
  This is a definition.
\end{proof}

\subsection{Picard-scheme predicates}
\label{subsec:fga_pic_scheme_carriers}

\begin{definition}[Existence of the Picard scheme]\leanok
  \label{def:has_pic_scheme}
  \dcref{custom}
  \lean{AlgebraicGeometry.Scheme.HasPicScheme}
  \uses{def:pic_sharp_carrier}
  The predicate \(\mathrm{HasPicScheme}(C)\) asserts that there exists a \(k\)-scheme
  \(X\), \emph{locally of finite type over \(k\)}, carrying a
  representability witness for the relative Picard functor
  \(\Pic^\sharp_{C/k}\) of \cref{def:pic_sharp_carrier}; in other words,
  that the Picard scheme \(\Pic_{C/k}\) exists and is locally of finite
  type. It is the existence assertion from which \cref{def:pic_scheme}
  extracts the representing object. The local-finiteness clause belongs to
  the same existence package as representability: Kleiman \S 4,
  Thm.~th:main exhibits \(\Pic_{C/k}\) as a disjoint union of open
  quasi-projective \(k\)-subschemes.

  This predicate concerns the \emph{unsheafified} functor, so its only witness
  \cref{def:inst_has_pic_scheme} is conditional on the \(k\)-rational point of
  \cref{def:has_rational_point}. Without a section the predicate is not available
  in general: Kleiman's pointless real conic (see the chapter introduction) is a
  curve satisfying these very hypotheses for which \(\Pic_{X/\bb R}\) is not
  representable. The unconditional statement, for the \'etale sheafification, is
  \cref{thm:fga_pic_representability_etale}.
  % Reason corrected 2026-07-30 (review-ajc): this said "since without a section
  % Pic^sharp is not a sheaf and the predicate is false in general". Both halves
  % were withdrawn earlier at the chapter introduction and missed here. The
  % Zariski-sheaf reason is about the ABSOLUTE Pic_X, not the relative functor
  % (which is Zariski-separated on these binders, th:cmp part 1); and the
  % non-representability is a CONDITIONAL Lean theorem
  % (not_exists_representing_picSharp_of_not_isIso) whose antecedent is quoted
  % from Kleiman rather than formalised, so "false in general" overstates it.
\end{definition}

\begin{proof}
  This is the defining existence predicate.
\end{proof}

\begin{definition}[\'Etale-sheafified relative Picard functor]\leanok
  \label{def:pic_et_sheaf}
  \dcref{custom}
  \lean{AlgebraicGeometry.Scheme.PicScheme.picEt}
  \uses{def:rel_pic_sharp}
  The associated \'etale sheaf \(\Pic_{(C/k)\et}\) of the relative Picard
  presheaf \(\Pic^\sharp_{C/k}\) of \cref{def:rel_pic_sharp}, on the site
  \((\Sch/k)\) equipped with the \'etale topology, with its abelian-group
  structure retained and (for representability) its underlying set-valued
  functor taken. This is Kleiman's \(\Pic_{(X/S)\et}\) of \S 2 (\texttt{df:Pfs}).

  The \'etale topology on \((\Sch/k)\) is the big \'etale topology on schemes
  localised at \(\Spec k\) in the sense of SGA 4 III 5.2.1: a sieve on a
  \(k\)-scheme \(T\) covers exactly when the sieve it generates on the
  underlying scheme covers \'etale-locally.
\end{definition}

\begin{proof}\leanok
  \uses{def:rel_pic_sharp}
  Sheafification of the group-valued presheaf of \cref{def:rel_pic_sharp} for
  the \'etale topology.
\end{proof}

\begin{lemma}[\(\Pic_{(C/k)\et}\) is a sheaf]\leanok
  \label{lem:pic_et_is_sheaf}
  \dcref{custom}
  \lean{AlgebraicGeometry.Scheme.PicScheme.picEt_isSheaf_forget}
  \uses{def:pic_et_sheaf}
  \(\Pic_{(C/k)\et}\) satisfies the sheaf axiom for the \'etale topology on
  \((\Sch/k)\), and so does its underlying set-valued functor.
\end{lemma}

\begin{proof}\leanok
  \uses{def:pic_et_sheaf}
  By construction: \cref{def:pic_et_sheaf} is a sheafification. The set-valued
  claim follows because the forgetful functor \(\Ab \to \Set\) preserves the
  limits the sheaf condition is expressed by.

  This lemma is what makes \cref{thm:fga_pic_representability_etale} askable with
  no hypothesis on \(C(k)\): the \'etale sheaf condition holds \emph{by
  construction} here, whereas for the unsheafified \(\Pic^\sharp_{C/k}\) Kleiman's
  pointless real conic is a curve on these very hypotheses where representability
  fails (chapter introduction).
  % Reason corrected 2026-07-30 (review-ajc): this said the unsheafified functor
  % "is not a sheaf", the withdrawn claim -- it is Zariski-SEPARATED on these
  % binders (th:cmp part 1); the absolute Pic_X is what fails. And the
  % non-representability is conditional in Lean
  % (not_exists_representing_picSharp_of_not_isIso), its antecedent quoted from
  % Kleiman, so "could not hold" is stronger than what is proved here.
\end{proof}

\begin{theorem}[FGA representability, \'etale formulation]
  \label{thm:fga_pic_representability_etale}
  \dcref{custom}
  \lean{AlgebraicGeometry.Scheme.fgaPicardRepresentability}
  \uses{def:pic_et_sheaf, lem:pic_et_is_sheaf, def:rel_pic_sharp}
  Let \(k\) be a field and \(C/k\) a smooth proper geometrically integral
  curve --- with \emph{no} hypothesis on \(C(k)\). Then:
  \begin{enumerate}
    \item the \'etale-sheafified relative Picard functor \(\Pic_{(C/k)\et}\) of
      \cref{def:pic_et_sheaf} is representable by a \(k\)-scheme that is
      separated and locally of finite type over \(k\); and
    \item if \(C\) has a \(k\)-rational point, the comparison
      \(\Pic^\sharp_{C/k} \to \Pic_{(C/k)\et}\) is an isomorphism.
  \end{enumerate}

  This is the project's central open obligation, and it is expected to remain
  open. It is the only assertion of this chapter that is not proved, and the
  Jacobian headline (\cref{chap:Jacobian}) rests on it.
\end{theorem}

\begin{proof}
  \uses{thm:rigid_pushforward_locally_free, thm:galois_quotient_points,
    lem:smooth_proper_quotient}
  Clause (1) is Kleiman \S 4, Thm.~\texttt{th:main} together with Cor.\
  \texttt{cor:algsch}: the Abel map exhibits the sheafified Picard functor as a
  quotient of a quasi-projective scheme by a smooth proper equivalence
  relation, and Altman--Kleiman descent (\cref{lem:smooth_proper_quotient})
  represents the quotient as a disjoint union of open quasi-projective
  \(k\)-subschemes. Along the Milne--Koll\'ar route the degree pieces
  \(\Pic^r\) are instead glued from the loci \(J^\Sigma\) over a separably
  closed field, descended to \(k\) by a finite Galois quotient, and assembled
  over \(\mathbb{Z}\) as a coproduct, which is locally of finite type because
  each piece is.

  Clause (2) is Kleiman \S 2, Thm.~2.5, as displayed in
  \cref{sec:fga_pic_setup}.

  The two clauses are stated together, rather than as separate assertions,
  because both are theorems of the same paper and neither is formalised, while
  the \(\Pic^\sharp\)-shaped statements later in this chapter need clause (2)
  to exist at all. Splitting clause (2) out would report the Jacobian headline
  as resting on six obligations where the mathematics has five.
\end{proof}

\begin{definition}[Picard scheme of a pointed curve]\leanok
  \label{def:inst_has_pic_scheme}
  \dcref{custom}
  \lean{AlgebraicGeometry.Scheme.picSchemeOfHasRationalPoint}
  \uses{def:has_pic_scheme, def:has_rational_point,
    thm:fga_pic_representability_etale}
  Let \(C/k\) be a smooth proper geometrically integral curve \emph{with a
  \(k\)-rational point}. Then the plain relative Picard functor
  \(\Pic^\sharp_{C/k}\) is represented by a scheme separated and locally of
  finite type over \(k\), so \cref{def:has_pic_scheme} holds.

  This is a \emph{conditional milestone}, strictly weaker than
  \cref{thm:fga_pic_representability_etale}, and it may not be presented as the
  project's headline. It is retained because it is true, because it records
  exactly what a rational point buys, and because the tangent-space development
  is written against the \(\Pic^\sharp\) shape.
\end{definition}

\begin{proof}\leanok
  \uses{thm:fga_pic_representability_etale}
  Both inputs are clauses of \cref{thm:fga_pic_representability_etale}: clause
  (1) supplies a scheme representing \(\Pic_{(C/k)\et}\) with the required
  properties, and clause (2) --- available because \(C\) has a rational point
  --- identifies \(\Pic^\sharp_{C/k}\) with \(\Pic_{(C/k)\et}\), so the same
  scheme represents it. Transport the representability witness along that
  isomorphism.
\end{proof}

\begin{remark}[What the representability statements of this chapter assert]
  \label{rem:representability_is_conditional}
  \dcref{custom}
  \uses{def:has_pic_scheme, def:inst_has_pic_scheme, def:pic_scheme,
    thm:fga_pic_representability}
  Every representability assertion of this chapter is an assertion \emph{about} the
  existence predicate \(\mathrm{HasPicScheme}\) of \cref{def:has_pic_scheme}, not a
  construction of the Picard scheme. \Cref{def:pic_scheme} extracts a scheme from that
  predicate, \cref{thm:fga_pic_representability} extracts the representing property, and
  \cref{thm:pic_is_group_scheme} equips the extracted object with its group structure; each
  is a theorem about what follows once the predicate holds.

  \emph{Corrected} (2026-07-29, \texttt{review-ajc}). This paragraph named
  \cref{def:inst_has_pic_scheme} as ``the one statement of this chapter that is asserted
  rather than proved''. That is no longer where the assertion lives, and the chapter
  introduction repeats the same misattribution. \Cref{def:inst_has_pic_scheme} is
  \emph{proved} --- its Lean witness \texttt{picSchemeOfHasRationalPoint} derives the
  \(\Pic^\sharp\) predicate from the two clauses of
  \cref{thm:fga_pic_representability_etale} plus the rational point, which is why it
  carries \texttt{\textbackslash leanok}. The single unproved statement of the chapter is
  \cref{thm:fga_pic_representability_etale} itself, a bare \texttt{sorry} in Lean
  (the Lean declaration \texttt{AlgebraicGeometry.Scheme.fgaPicardRepresentability});
  the \(\Pic^\sharp\) predicate is one of
  its consequences, not its source. The older reading dates from when the predicate was
  supplied by an assumed instance, which was deleted on 2026-07-28 under the owner decision
  recorded at \cref{rem:rational_point_scope}.

  The distinction is easy to lose because a statement quantifying over the predicate is
  unconditional as a statement --- its hypothesis is discharged by whoever applies it. It
  becomes visible one step later, at a curve where the predicate must be supplied: there
  every consequence above depends on \cref{def:inst_has_pic_scheme}. So the correct reading
  of this chapter is that the passage \emph{from} representability to the Picard scheme and
  its group structure is complete, and representability is not.

  Two things sharpen this since the \'etale formulation was adopted. First, the
  \(\Pic^\sharp\)-shaped predicate is supplied only for a \emph{pointed} curve, so the
  consequences above are conditional statements in a second and independent way: their
  hypothesis is not just unproved but unavailable over a general field. Second, the
  corresponding \'etale predicate \emph{is} supplied unconditionally, and everything said
  here applies to it verbatim with
  \cref{thm:fga_pic_representability_etale} in place of
  \cref{def:inst_has_pic_scheme} --- the passage from representability to the Picard
  scheme, its group structure, its separatedness and its local finiteness is complete
  there too, and representability is still the one thing that is not.
\end{remark}

\begin{definition}[Representability of \(\Pic^\sharp_{C/k}\)]\leanok
  \label{def:pic_sharp_representable}
  \dcref{custom}
  \lean{AlgebraicGeometry.Scheme.PicScheme.PicSharpRepresentable}
  \uses{def:pic_sharp_carrier, def:pic_scheme}
  The predicate \(\mathrm{PicSharpRepresentable}(C)\) asserts that the specific scheme
  \(\Pic_{C/k}\) of \cref{def:pic_scheme} is a representing object for
  \(\Pic^\sharp_{C/k}\). It is the assertion from which the representability
  witness \cref{thm:fga_pic_representability} is extracted.
\end{definition}

\begin{proof}
  This is a definition.
\end{proof}

\begin{definition}[Representability identification]\leanok
  \label{def:inst_pic_sharp_representable}
  \dcref{custom}
  \lean{AlgebraicGeometry.Scheme.PicScheme.instPicSharpRepresentable}
  \uses{def:pic_sharp_representable, def:has_pic_scheme}
  The scheme \(\Pic_{C/k}\) is the chosen witness of
  \cref{def:has_pic_scheme}; the first component of that witness identifies
  it as a representing object for \(\Pic^\sharp_{C/k}\).

  \emph{Scope} (added 2026-07-29, \texttt{review-ajc}). This is a projection of
  the hypothesis \cref{def:has_pic_scheme}, not an independent result: its
  \texttt{\textbackslash leanok} certifies only that extraction from the
  existence package elaborates. And that hypothesis is currently
  \emph{uninhabited} in Lean --- \(\mathrm{HasPicScheme}\) has no instance, its
  only producer requires a \(k\)-rational point, and no unconditional producer of
  such a point exists either. Every statement in this section quantifying over
  \cref{def:has_pic_scheme} is therefore conditional on an assumption nothing
  can presently supply; the unconditional interface is the \'etale one.
\end{definition}

\begin{proof}\leanok
  \uses{def:has_pic_scheme}
  This is the representing component of the existence statement in
  \cref{def:has_pic_scheme}.
\end{proof}

\begin{definition}[Local finite type of \(\Pic_{C/k}\)]\leanok
  \label{def:pic_scheme_lft}
  \dcref{custom}
  \lean{AlgebraicGeometry.Scheme.PicScheme.PicSchemeLocallyOfFiniteType}
  \uses{def:pic_scheme, def:has_pic_scheme}
  The predicate
  \[
    \mathrm{PicSchemeLocallyOfFiniteType}(C)
  \]
  asserts that the
  structural morphism of the Picard scheme \(\Pic_{C/k}\) of
  \cref{def:pic_scheme} is locally of finite type --- part (1) of Kleiman
  \S 4, Thm.~th:main (\(\Pic_{C/k}\) is a disjoint union of open
  quasi-projective \(k\)-subschemes). It is the hypothesis the
  identity-component construction of
  \cref{chap:Picard_IdentityComponent} uses to specialise
  \(G^0\) to \(G = \Pic_{C/k}\). It follows from the local-finiteness
  component of \cref{def:has_pic_scheme}.
\end{definition}

\begin{proof}\leanok
  This is a definition.
\end{proof}

\begin{definition}[Local-finiteness assertion]\leanok
  \label{def:inst_pic_scheme_lft}
  \dcref{custom}
  \lean{AlgebraicGeometry.Scheme.PicScheme.instPicSchemeLocallyOfFiniteType}
  \uses{def:pic_scheme_lft, def:has_pic_scheme}
  Whenever \cref{def:has_pic_scheme} holds, the second component of its
  chosen witness gives local finiteness of \(\Pic_{C/k}\). Thus
  representability and local finiteness are supplied by one existence
  package, as in Kleiman \S 4, Thm.~th:main.
\end{definition}

\begin{proof}\leanok
  \uses{def:has_pic_scheme}
  This is the local-finiteness component of the existence statement in
  \cref{def:has_pic_scheme}.
\end{proof}

\section{Dependencies of the representability theorem}
\label{sec:fga_pic_sorry_closure_order}

The construction of \(\Pic_{C/k}\) separates into three mathematical inputs:
the relative-divisor functor, its Abel map to the relative Picard functor, and
the FGA existence theorem. The first two supply concrete functors and a
natural transformation. The third combines them with regularity, cohomology
and base change, and effective descent to produce a scheme locally of finite
type representing \(\Pic^\sharp_{C/k}\).

\subsection{The Picard-scheme existence theorem}
\label{subsec:sorry_has_pic_scheme}

For a
smooth proper geometrically integral curve \(C/k\) with a \(k\)-rational
point, the relative Picard functor \(\Pic^\sharp_{C/k}\)
(\cref{def:pic_sharp_carrier}) is representable by a \(k\)-scheme. Either route
of \cref{sec:fga_pic_setup} establishes it, and they consume different inputs.

Along the Milne--Koll\'ar route of \cref{sec:fga_pic_milne_kollar} --- the one
being formalised --- the inputs are: the rigidified pushforward
\cref{thm:rigid_pushforward_gate}, which holds for every curve satisfying this
chapter's hypotheses and supplies both the canonical divisor family and the
semicontinuity of \(h^0\) cutting out the loci \(C^\Sigma\); uniform \(H^1\) vanishing in high degree, which makes those loci
cover; representability of the degree-\(r\) divisor functor, which supplies the
ambient \(Z\) together with its Grassmannian immersion; and the finite Galois
quotient \cref{thm:galois_quotient_points} with
\cref{lem:stable_affine_open}, which descends the result from a separably closed
field to \(k\). No Hilbert or Quot scheme occurs among them.

Along the quotient route the proof first applies Kleiman \S 2, Thm.~2.5, using
the section to identify the plain functor with its \'etale sheafification, and
then uses Kleiman's four steps
(\cref{thm:fga_pic_representability}): Hilbert-polynomial stratification,
\(m\)-regularity bound, Abel-map factorisation through
\(\mathrm{Div}_{C/k}\), and the smooth-proper quotient
(\cref{lem:smooth_proper_quotient}). Its inputs are the Quot/Hilbert
construction of \cref{chap:Picard_QuotScheme}, representability of
\(\mathrm{Div}_{C/k}\) as in \cref{rem:div_functor_representability}, the
regularity result \cref{cor:flattening_stratification_curve}, the base-change
result \cref{lem:cech_flat_base_change}, and Altman--Kleiman effective descent.
Both routes rest on the same divisor and comparison substrate. A section is
used only for the pointed specialisation, where it identifies the plain
relative Picard functor with its \'etale sheafification. The arbitrary-field
\'etale statement has no rational-point hypothesis: in the Milne--Koll\'ar
construction, rational points are introduced only after base change, and the
resulting \'etale Picard object is descended to \(k\).

Either way, the resulting representing object gives
natural bijections
\[
  (T\to\Pic_{C/k})\simeq
  \Pic(C\times_kT)/\pi_T^*\Pic(T),
\]
including the dual-number case \(T=\Spec k[\epsilon]\) used in the
tangent-space calculation.

\subsection{The relative-divisor input}
\label{subsec:sorry_has_div_functor}

The relative-divisor functor \(\Div_{X/S}\) of
\cref{def:div_functor} is the subfunctor of the Hilbert functor
\(\Hilb_{C/k} = \Quot_{\mathcal{O}_C/C/k}\) of relative effective Cartier
divisors, in the invertible-kernel quotient encoding of
\cref{def:div_family}. The presheaf \(\mathrm{Div}_{C/k}\) of
\cref{def:div_functor_carrier} is its instantiation at
\(C\to\Spec k\). Its divisor condition is stable under base change by
\cref{lem:relative_divisor_base_change}, while representability of
\(\Div_{C/k}\) by an open subscheme of the Hilbert scheme (Kleiman \S 3,
Thm.~th:repDiv) is the content of
\cref{rem:div_functor_representability}.

\subsection{The Abel-map input}
\label{subsec:sorry_has_abel_map}

The class \cref{def:has_abel_map} carries the natural transformation whose
value on a divisor family is
\[
  \mathrm{abelMapWitness}(C)\langle\mathcal F,q\rangle=-[\ker q].
\]
It is the composite of the kernel-class natural transformation
\(\mathrm{abelKernelNatTrans}(C)\) (naturality via the kernel--flat-base-change
comparison isomorphism) with the group-presheaf negation
\(\mathrm{picNeg}(C)\). Thus the Abel map has the defining property in
\cref{lem:line_bundle_quot_correspondence}; it is a natural transformation
of the presheaves in \cref{def:div_functor_carrier} and
\cref{def:pic_sharp_carrier}.

\subsection{Assembly of the inputs}
\label{subsec:fga_pic_closure_order_summary}

The FGA construction consumes the relative-divisor functor and Abel map,
then applies the regularity, base-change, and effective-quotient results of
\cref{thm:fga_pic_representability}. For the pointed specialisation, a map
\(P:\mathbf 1\to C\) identifies the plain relative Picard functor with its
sheafification. The arbitrary-field statement instead represents the \'etale
sheafification directly and has no rational-point hypothesis. The resulting
representability and local-finiteness statements feed the group-scheme and
Jacobian constructions.

\section{Related constructions}
\label{sec:fga_pic_out_of_scope}

The representable Picard scheme and its group structure interact with the
following constructions, treated in the indicated companion chapters.
\begin{itemize}
  \item \textbf{Identity component \(\Pic^0_{C/k}\)} --- the open and
    closed subgroup scheme of \(\Pic_{C/k}\) containing the unit and
    consisting of those classes lying in the connected component of
    \(\Pic_{C/k}\). Its dimension, smoothness, and abelian-variety
    structure are treated in \cref{chap:Picard_IdentityComponent} and
    \cref{chap:Picard_Pic0AbelianVariety}.
  \item \textbf{Quot scheme representability} --- consumed by
    \cref{thm:fga_pic_representability} via \texttt{thm:quot\_representable}
    (\cref{chap:Picard_QuotScheme}), not re-proved here. It is an input of the
    quotient route only; the Milne--Koll\'ar route of
    \cref{sec:fga_pic_milne_kollar} does not use it, although it does use the
    Grassmannian and flattening substrate of the same chapters.
  \item \textbf{\'Etale sheafification.} The relative Picard presheaf and its
    topology-parametric sheafification are constructed in
    \cref{chap:Picard_RelPicFunctor}. Under the
    standing rational-point hypothesis, Kleiman \S 2, Thm.~2.5 identifies
    the plain relative functor with the \'etale sheafification.
  \item \textbf{Universal Poincar\'e sheaf} (Kleiman~\S 4, Ex.~4.3)
    --- exists when \(\Pic_{C/k}\) represents \(\Pic_{C/k}\) itself and
    \(C/k\) has a section. Under the chapter's standing rational-point
    hypothesis both conditions in fact hold (Thm.~2.5 identifies the
    plain functor with the represented sheaf), so the Poincar\'e sheaf
    is then available as an associated construction.
  \item \textbf{Mumford's example} (Kleiman~\S 4, Eg.~4.14), counter-examples
    to representability for non-integral geometric fibers --- not
    relevant to a smooth proper geometrically integral curve.
\end{itemize}

\section{Dependency structure}
\label{sec:fga_pic_consistency_check}

The principal constructions form a connected dependency chain rooted in
\cref{chap:Picard_LineBundlePullback}, \cref{chap:Picard_QuotScheme}, and
\cref{chap:Picard_RelPicFunctor}:
\begin{itemize}
  \item \cref{thm:rigid_pushforward_locally_free} uses \cref{lem:p1_integral},
    and \cref{lem:stable_affine_open} uses
    \cref{thm:galois_quotient_points}, which uses
    \cref{thm:speiser_descent}. These are the inputs specific to the
    Milne--Koll\'ar route (\cref{sec:fga_pic_milne_kollar}).
  \item \cref{lem:line_bundle_quot_correspondence} uses
    \cref{def:line_bundle_on_product} and
    \cref{thm:pullback_natural}.
  \item \cref{lem:smooth_proper_quotient} is the structural quotient
    statement used by the quotient-route representability argument.
  \item \cref{thm:fga_pic_representability} uses
    \cref{lem:line_bundle_quot_correspondence},
    \texttt{thm:quot\_representable},
    \texttt{def:rel\_pic\_sharp},
    \cref{def:has_rational_point},
    \cref{thm:relative_pic_quotient_well_defined}, and
    \cref{lem:smooth_proper_quotient}.
  \item \cref{thm:pic_is_group_scheme} uses
    \cref{def:pic_scheme}, \cref{thm:fga_pic_representability},
    \texttt{def:rel\_pic\_sharp}, and
    \cref{thm:relative_pic_quotient_well_defined}.
  \item \cref{def:pic_scheme} uses
    \cref{def:has_pic_scheme} and \cref{def:inst_has_pic_scheme}.
\end{itemize}
The Quot-scheme theorem used above is proved in
\cref{chap:Picard_QuotScheme}, while the relative Picard functor is defined in
\cref{chap:Picard_RelPicFunctor}.
